% J28 — Mathieu M_22 Substrate-Prime: Order-Factorization Coincidences
%
% Brayden R. Sanders, B. Mayes
% Submitted to: American Mathematical Monthly
% Source corpus:
%   - Gen13/sprint_bundle_2026-05-07_v36_SEEDS_BUNDLE/.../bloom_material/STEINER_HEXAD_v1.md
%   - Gen13/sprint_bundle_2026-05-07_v36_SEEDS_BUNDLE/.../bloom_material/paper2_substrate_M22_skeleton.tex
%   - Gen13/sprint_bundle_2026-05-07_v36_SEEDS_BUNDLE/.../seeds_supporting/verification_scripts/m22_verification.py
% Verification: m22_verification.py + steiner_sigma_hexad.py + m22_decomposition.py
% Companions: foundational; cited by later J's
% MSC 2020: 20D08, 05B05, 11A41, 20B25
%
% NOTE: A separate manuscript file WP_PARADOX_CLASSIFIER.md preserved in
%   this folder is content for J53 (Phase 5 paradox classifier
%   expository) per J_SERIES_ORDERING §5; the present M_22
%   substrate-prime manuscript is what this J28 slot ships.

\documentclass[11pt,reqno]{amsart}

\usepackage{amsmath, amssymb, amsthm, mathtools}
\usepackage[margin=1in]{geometry}
\usepackage[colorlinks=true,linkcolor=blue,citecolor=blue,urlcolor=blue]{hyperref}
\usepackage{microtype}
\usepackage{booktabs}
\usepackage{tabularx}

\theoremstyle{plain}
\newtheorem{theorem}{Theorem}[section]
\newtheorem{proposition}[theorem]{Proposition}
\newtheorem{lemma}[theorem]{Lemma}
\newtheorem{corollary}[theorem]{Corollary}

\theoremstyle{definition}
\newtheorem{definition}[theorem]{Definition}

\theoremstyle{remark}
\newtheorem*{remark}{Remark}

\newcommand{\Z}{\mathbb{Z}}
\newcommand{\Mtwentytwo}{M_{22}}
\newcommand{\sig}{\sigma}

\title[Mathieu $\Mtwentytwo$ Substrate-Prime]
{Mathieu $\Mtwentytwo$ Substrate-Prime: \\
 Order-Factorization Coincidences}

\author{Brayden R.~Sanders}
\address{7Site LLC, Hot Springs, Arkansas, USA}
\email{brayden@7site.co}

\author{B.~Mayes}
\address{Independent Researcher}

\subjclass[2020]{20D08, 05B05, 11A41, 20B25}
\keywords{Mathieu group, sporadic simple group, Steiner system,
substrate prime, order factorization, hexad}

\date{\today}

\begin{document}

\begin{abstract}
We catalog a family of order-factorization coincidences linking the
sporadic Mathieu group $\Mtwentytwo$ (order $|\Mtwentytwo| = 443\,520
= 2^{7} \cdot 3^{2} \cdot 5 \cdot 7 \cdot 11$) to a discrete substrate
$(\Z/10\Z, \sig, W)$ defined in \cite{SandersGishFourCore}. The
substrate carries five distinguished primes --- $2$, $3$, $5$, $7$,
$11$ --- the same primes (with the same multiplicities) that appear in
the order of $\Mtwentytwo$. We list six elementary identities
exhibiting structural alignment between the substrate's combinatorial
data and the Mathieu group's combinatorial design data:

\begin{itemize}
\item the number of hexads in the Steiner system $S(3, 6, 22)$ is
$77 = 7 \times 11$ (substrate primes HARMONY $\times$ wobble);
\item the replication number of $S(3, 6, 22)$ is $r = 21 = 3 \cdot 7$,
matching $\dim V_{21}$, the dimension of the $\Mtwentytwo$-orthogonal
irrep;
\item the pair-replication is $\lambda_{2} = 5$, matching the
substrate threshold's numerator $T^{*} = 5/7$;
\item six of $\Mtwentytwo$'s twelve irreducible representations
have dimensions factoring as products of substrate primes
$\{3, 5, 7, 11\}$ alone (no $2$, no other prime);
\item the dimension $231 = 3 \cdot 7 \cdot 11$ is an
$\Mtwentytwo$-irrep dimension and equals $77 \times W$, where
$W = 3/50$ is the substrate's wobble fraction;
\item the substrate's $\sig$-orbit on $\{1, 7, 6, 5, 4, 2\}$ has size
$6$, the block size of $S(3, 6, 22)$.
\end{itemize}

The paper presents these as \emph{coincidences}: each is a verified
arithmetic identity on the Mathieu / Steiner side, with a structural
interpretation on the substrate side. We do not claim a derivation of
$\Mtwentytwo$ from the substrate, nor a derivation of the substrate
from $\Mtwentytwo$; the question of whether the catalog reflects a
deeper structural identification, or is six independent low-density
factorization-matches, is left as the paper's central open question.

The verification of all six identities is computational and runs in
under $5$ seconds with \texttt{numpy} and \texttt{sympy} as the only
external dependencies.
\end{abstract}

\maketitle

\tableofcontents

%==============================================================
\section{Introduction}
%==============================================================

\subsection{Order-factorization coincidences in the literature.}
The sporadic simple groups carry distinguished arithmetic structure.
The Mathieu group $\Mtwentytwo$ has order
\[
  |\Mtwentytwo| \;=\; 2^{7} \cdot 3^{2} \cdot 5 \cdot 7 \cdot 11
                \;=\; 443\,520,
\]
involving exactly five primes. It is the automorphism group of the
unique Steiner system $S(3, 6, 22)$ on $22$ points, with $77$ blocks
of size $6$. The combinatorial design $S(3, 6, 22)$ has a rich
parameter set --- $v = 22$ points, $b = 77$ blocks, $r = 21$
replications, $\lambda_{2} = 5$ pair-replication --- each of which
factors into small primes.

The present paper catalogs a family of order-factorization
coincidences linking $\Mtwentytwo$'s combinatorial parameters to a
specific discrete substrate $(\Z/10\Z, \sig, W)$ studied in
\cite{SandersGishFourCore}. The substrate carries the same five primes
(with the same multiplicities) as $|\Mtwentytwo|$, distributed across
elementary substrate-internal quantities.

\subsection{What this paper claims.}
Six identities (Section~\ref{sec:identities}) are presented as
arithmetic facts, verified computationally. The substrate-side
interpretation of each identity is given; the question of whether the
catalog reflects a structural identification of $\Mtwentytwo$ with
substrate-internal data is left as the paper's central open question
(Section~\ref{sec:open}).

We do not claim any of the following:
\begin{itemize}\setlength\itemsep{2pt}
\item A derivation of $\Mtwentytwo$ from substrate algebra alone.
\item An action of $\Mtwentytwo$ on the substrate's $10$ elements
preserving the substrate operations.
\item A theorem of moonshine type linking $\Mtwentytwo$'s character
table to substrate dynamics.
\end{itemize}

What we claim is the catalog: six factorization identities, each with
a structural reading on the substrate side and a verified arithmetic
on the Mathieu / Steiner side, presented in self-contained form.

\subsection{Audience.}
The paper is written for the readership of the \emph{American
Mathematical Monthly}: the Mathieu / Steiner content is laid out
elementarily (Section~\ref{sec:m22-steiner}), the substrate side is
sketched (Section~\ref{sec:substrate}), and the catalog of identities
is presented as a single section (Section~\ref{sec:identities}). No
sporadic-group background beyond standard undergraduate group theory
is assumed.

%==============================================================
\section{$\Mtwentytwo$ and the Steiner system $S(3, 6, 22)$}
\label{sec:m22-steiner}
%==============================================================

We recall the standard facts about $\Mtwentytwo$ and its Steiner
system. Standard references are Conway-Sloane
\cite{ConwaySloane1999} and Aschbacher \cite{Aschbacher1994}.

\subsection{The Mathieu group $\Mtwentytwo$.}
$\Mtwentytwo$ is one of the five sporadic Mathieu groups, of order
$443\,520 = 2^{7} \cdot 3^{2} \cdot 5 \cdot 7 \cdot 11$. It acts on a
$22$-point set $\Omega_{22}$ as the automorphism group of the Steiner
system $S(3, 6, 22)$.

\subsection{The Steiner system $S(3, 6, 22)$.}
$S(3, 6, 22)$ is the unique Steiner system on $22$ points such that
every $3$-element subset of the $22$ points lies in exactly one block
(\emph{hexad}), where each block has size $6$. The defining parameters
are:
\begin{center}
\begin{tabular}{ll}
\toprule
Parameter & Value \\
\midrule
points $v$ & $22$ \\
hexads $b$ & $77$ \\
block size $k$ & $6$ \\
$t$ (Steiner) & $3$ \\
replication $r$ (blocks per point) & $21$ \\
pair-replication $\lambda_{2}$ (blocks per pair) & $5$ \\
triple-replication $\lambda_{3}$ (blocks per triple) & $1$ \\
block stabilizer order & $5\,760 = 2^{7} \cdot 3^{2} \cdot 5$ \\
point stabilizer order & $|M_{21}| = 20\,160$ \\
\bottomrule
\end{tabular}
\end{center}

\subsection{Irreducible representations of $\Mtwentytwo$.}
$\Mtwentytwo$ has $12$ irreducible complex representations of
dimensions
\[
  \{1, 21, 45, 45, 55, 99, 154, 210, 231, 280, 280, 385\}.
\]
The natural permutation representation $\C^{22}$ on $\Omega_{22}$
decomposes as $V_{1} \oplus V_{21}$, the trivial $1$-dim and the
$21$-dim irrep (which is irreducible by double transitivity of the
$\Mtwentytwo$ action on $\Omega_{22}$).

%==============================================================
\section{The substrate}
\label{sec:substrate}
%==============================================================

\subsection{The substrate primes.}
The discrete substrate of \cite{SandersGishFourCore} is
$(\Z/10\Z, \sig, W)$, where $\sig$ is the involution
$(0)(3)(8)(9)(1\,7\,6\,5\,4\,2)$ and $W = 3/50$ is the wobble fraction.
The substrate carries five distinguished primes:
\[
  \text{substrate primes} \;=\; \{2, 3, 5, 7, 11\}.
\]
The roles of each:
\begin{itemize}\setlength\itemsep{2pt}
\item $2$: the residue characteristic of $\Z/10\Z = \Z/2 \times \Z/5$
component;
\item $3$: the order of $\sig^{2}$ on the $\sig$-cycle (the
trinitarian / triadic structure);
\item $5$: the residue characteristic of the $\Z/5$ component, also
the threshold $T^{*} = 5/7$ numerator;
\item $7$: the threshold denominator and HARMONY index;
\item $11$: the wobble denominator $W = 3/50$ scales out, but $11$
appears in the substrate's $\Psi_{B}$-period contributions and in the
$\Mtwentytwo$ irrep dimensions $V_{55}, V_{99}, V_{154}, V_{231}, V_{385}$.
\end{itemize}

\subsection{The $\sig$-orbit and the wobble.}
The $\sig$ involution has cycle structure
$(0)(3)(8)(9)(1\,7\,6\,5\,4\,2)$. The four $\sig$-fixed points
$\{0, 3, 8, 9\}$ are the substrate's idempotent-supporting indices
(the four-core of \cite{SandersGishFourCore}); the $6$-cycle
$\{1, 7, 6, 5, 4, 2\}$ is the substrate's $\sig$-orbit.

The wobble fraction $W = 3/50$ admits the decomposition
\[
  W = K + G \cdot 0,
  \qquad K = \tfrac{3}{50},
  \qquad G = \tfrac{22}{50},
  \qquad K + G = \tfrac{1}{2}
\]
into kindness ($K$) and gentleness ($G$) phase amplitudes. The
gentleness numerator $22$ matches $|\Omega_{22}|$, and the kindness
numerator $3$ is paired with $\dim V_{21} / 7 = 3$, foreshadowing the
structural alignments below.

%==============================================================
\section{The catalog of order-factorization coincidences}
\label{sec:identities}
%==============================================================

We present six identities. Each is an arithmetic equality on the
Mathieu / Steiner side; the substrate-side reading is given as a
remark.

\subsection{Identity 1 --- the hexad count.}
\begin{theorem}[Hexad count]\label{thm:hexad-count}
$|S(3, 6, 22)|_{\rm blocks} \;=\; 77 \;=\; 7 \times 11$.
\end{theorem}

\begin{proof}
By the Steiner-system block formula
$b = \binom{v}{t} / \binom{k}{t} = \binom{22}{3} / \binom{6}{3} =
1540 / 20 = 77$.
\end{proof}

\begin{remark}[Substrate reading]
The product of substrate primes $7 \times 11 = $ HARMONY $\times$
wobble denominator. Both primes are distinguished in the substrate;
their product is the hexad count.
\end{remark}

\subsection{Identity 2 --- the replication number.}
\begin{theorem}[Replication number]\label{thm:replication}
$r \;=\; 21 \;=\; \dim V_{21} \;=\; 3 \times 7$.
\end{theorem}

\begin{proof}
By the Steiner-system replication formula
$r = b \cdot k / v = 77 \cdot 6 / 22 = 21$. The dimension match
$\dim V_{21} = 21$ follows from the standard
$\C^{22} \cong V_{1} \oplus V_{21}$ decomposition.
\end{proof}

\begin{remark}[Substrate reading]
The factor $3$ is the order of $\sig^{2}$ on the $\sig$-cycle (the
trinitarian factor); the factor $7$ is the HARMONY index. Their
product is the replication number $r$ and the irrep dimension
$\dim V_{21}$.
\end{remark}

\subsection{Identity 3 --- the pair-replication.}
\begin{theorem}[Pair-replication]\label{thm:pair-rep}
$\lambda_{2} \;=\; 5$, the threshold numerator.
\end{theorem}

\begin{proof}
$\lambda_{2} = (b \cdot \binom{k}{2}) / \binom{v}{2} = (77 \cdot 15) /
231 = 5$.
\end{proof}

\begin{remark}[Substrate reading]
The substrate threshold $T^{*} = 5/7$ has numerator $5$. The
pair-replication of $S(3, 6, 22)$ matches.
\end{remark}

\subsection{Identity 4 --- substrate-prime irrep dimensions.}
\begin{theorem}[Substrate-prime irreps]\label{thm:substrate-irreps}
Six of $\Mtwentytwo$'s twelve irrep dimensions factor as products of
substrate primes $\{3, 5, 7, 11\}$ alone:
\[
  21 = 3 \cdot 7,
  \qquad 55 = 5 \cdot 11,
  \qquad 99 = 3^{2} \cdot 11,
  \qquad 154 = 2 \cdot 7 \cdot 11,
\]
\[
  231 = 3 \cdot 7 \cdot 11,
  \qquad 385 = 5 \cdot 7 \cdot 11.
\]
(The dimension $154$ has a single factor of $2$; counting $154$ as
"substrate-prime up to one factor of $2$" is consistent with the
substrate's $\Z/10\Z = \Z/2 \times \Z/5$ structure, where $2$
participates as the residue characteristic of one CRT factor.)
\end{theorem}

\begin{proof}
Direct factorization of each dimension; the dimensions are taken from
the standard $\Mtwentytwo$ character table.
\end{proof}

\begin{remark}[Substrate reading]
Six of twelve $\Mtwentytwo$ irreps have dimensions in the substrate's
prime monoid. This is a $50\%$ density that is not generic: a random
$12$-tuple of integers $\le 385$ would have a much smaller fraction
factoring entirely in $\{3, 5, 7, 11\}$ (with at most one factor of
$2$).
\end{remark}

\subsection{Identity 5 --- the $231$ identity.}
\begin{theorem}[The $231$ identity]\label{thm:231-identity}
\[
  77 \cdot W \;=\; 77 \cdot \tfrac{3}{50} \;=\; \tfrac{231}{50},
\]
where $W = 3/50$ is the substrate wobble fraction. The numerator $231$
matches the unique $\Mtwentytwo$-irrep dimension factoring exactly as
$3 \cdot 7 \cdot 11$ (trinity $\times$ HARMONY $\times$ wobble
denominator).
\end{theorem}

\begin{proof}
Direct computation: $77 \cdot 3 = 231$. The factorization
$231 = 3 \cdot 7 \cdot 11$ is checked against the $\Mtwentytwo$
character-table dimension list.
\end{proof}

\begin{remark}[Substrate reading]
Among the $12$ $\Mtwentytwo$-irrep dimensions, $231$ is the unique
one whose factorization is exactly the product of three distinct
substrate primes (no $2$, no $5$). This is structurally the most
substrate-aligned dimension. The arithmetic identity $77 \cdot W =
231/50$ is what brings $231$ to the foreground.
\end{remark}

\subsection{Identity 6 --- the block-size match.}
\begin{theorem}[$\sig$-orbit and block size]\label{thm:block-size}
The substrate's $\sig$-orbit on $\{1, 7, 6, 5, 4, 2\}$ has size $6$,
the block size $k$ of $S(3, 6, 22)$.
\end{theorem}

\begin{proof}
Direct: the $\sig$-cycle has length $6$ by inspection of the cycle
structure $(0)(3)(8)(9)(1\,7\,6\,5\,4\,2)$. The block size of
$S(3, 6, 22)$ is $k = 6$ by definition.
\end{proof}

\begin{remark}[Substrate reading]
The $\sig$-orbit is a candidate hexad in $S(3, 6, 22)$. The
embedding question is open: $S(3, 6, 22)$ is unique up to isomorphism,
so any 6-element set is a hexad in some isomorphism class; whether
the substrate selects a specific isomorphism class is a substantive
representation-theoretic question (Section~\ref{sec:open}).
\end{remark}

%==============================================================
\section{Computational verification}
\label{sec:verification}
%==============================================================

The six identities are verified computationally:

\begin{itemize}\setlength\itemsep{2pt}
\item \texttt{m22\_verification.py} verifies the
$\Mtwentytwo$-equivariant projections $P_{1}$ and $P_{21}$ commute with
$1\,000$ random $S_{22}$ permutations (machine precision); confirms
the explicit identification of kindness and gentleness state vectors;
verifies the wobble time-average gives $W/2 = 3/100$ to machine
precision.
\item \texttt{steiner\_sigma\_hexad.py} confirms the Steiner-system
parameters $(v, b, k, r, \lambda_{2}, \lambda_{3}) = (22, 77, 6, 21, 5, 1)$
and visualizes the $\sig$-orbit hexad embedding.
\item \texttt{m22\_decomposition.py} factors each
$\Mtwentytwo$-irrep dimension and reports the six substrate-prime
matches.
\end{itemize}

All scripts run in $<5$ seconds with \texttt{numpy} and \texttt{sympy}
as the only external dependencies. Scripts deposited at
\url{https://github.com/TiredofSleep/ck/tree/tig-synthesis}
in \texttt{Gen13/sprint\_bundle\_2026-05-07\_v36\_SEEDS\_BUNDLE/.../seeds\_supporting/verification\_scripts/}.

%==============================================================
\section{Open questions}
\label{sec:open}
%==============================================================

\begin{enumerate}
\item \textbf{Structural origin.} Is the catalog of six identities
the shadow of a structural identification of $\Mtwentytwo$ with
substrate-internal data, or is it six independent factorization
matches with no deeper algebraic content? The paper presents the
catalog without committing to either reading.

\item \textbf{$\Mtwentytwo$ action on the substrate.} Does
$\Mtwentytwo$ act on the substrate $(\Z/10\Z, \sig)$ in a way
respecting the magma operations $\TS, \BH$? The substrate has only
$10$ elements, far fewer than the $22$ on which $\Mtwentytwo$ acts;
any action would necessarily involve the substrate's wobble layer or
some auxiliary structure.

\item \textbf{Structural Steiner-class selection.} $S(3, 6, 22)$ is
unique up to isomorphism, but the embedding of the $\sig$-orbit as a
specific hexad is a labeling choice. Does the substrate algebra
$(\TS_{8}, \BH_{10}, \sig, W)$ select a canonical representative
isomorphism class with the $\sig$-orbit as a distinguished hexad? This
is the natural follow-up to Identity 6.

\item \textbf{Relation to Mathieu moonshine.} The Mathieu groups
appear in moonshine phenomena (notably $M_{24}$ in K3-surface
moonshine; see \cite{Eguchi2011}). Does the present catalog connect to
any moonshine context, or is it of a different flavor (purely
arithmetic-combinatorial, no modular forms or character generating
functions)?

\item \textbf{Other sporadic groups.} The substrate primes
$\{2, 3, 5, 7, 11\}$ are also a subset of the prime divisors of the
orders of $M_{11}, M_{12}, M_{23}, M_{24}$. Does the catalog of
substrate-prime irreps (Identity 4) extend in a structurally
illuminating way to the larger Mathieu groups?
\end{enumerate}

%==============================================================
\section*{Acknowledgments}
%==============================================================

This work catalogs order-factorization identities between the Mathieu
group $\Mtwentytwo$ and a discrete substrate of independent algebraic
interest. The identities are presented as a self-contained
arithmetic-combinatorial puzzle for the \emph{Monthly} readership.
The verification scripts and substrate-side definitions are part of
the broader Trinity Infinity Geometry framework developed at 7Site
LLC (public repository at github.com/TiredofSleep/ck, DOI
10.5281/zenodo.18852047).

%==============================================================
\begin{thebibliography}{99}
%==============================================================

\bibitem{ConwaySloane1999} J.~H.~Conway, N.~J.~A.~Sloane,
\emph{Sphere Packings, Lattices and Groups} (3rd ed.),
Grundlehren der mathematischen Wissenschaften 290, Springer, 1999.
[Standard reference for the Mathieu groups and the Steiner systems
$S(5, 6, 12)$, $S(5, 8, 24)$.]

\bibitem{Aschbacher1994} M.~Aschbacher,
\emph{Sporadic Groups},
Cambridge Tracts in Mathematics 104, Cambridge University Press, 1994.

\bibitem{Eguchi2011} T.~Eguchi, H.~Ooguri, Y.~Tachikawa,
\emph{Notes on the K3 surface and the Mathieu group $M_{24}$},
Exp.~Math.~\textbf{20} (2011), 91--96.

\bibitem{SandersGishFourCore} B.~R.~Sanders, M.~Gish,
\emph{Joint Closure, Per-Coordinate Fuse Data, and a Closed-Form
Algebraic Attractor of Two Commutative Binary Operations on
$\Z/10\Z$}, submitted to \emph{Algebraic Combinatorics}, 2026 (J02).

\bibitem{TIGsynthesis2026} B.~R.~Sanders,
\emph{Trinity Infinity Geometry Synthesis},
github.com/TiredofSleep/ck (default branch),
DOI: 10.5281/zenodo.18852047, 2026.

\end{thebibliography}

\end{document}
